% -------------------------------
% Plantilla del TFM en la ESIIAB/ESICR de la UCLM. Versión 0.1 (Preliminar)
%
% Luis de la Ossa
% 
% Compilar con XeLaTeX 
% -------------------------------


% -------------------------------
% Segunda portada
% -------------------------------
\cleardoublepage
\setcounter{page}{1} \null


\begin{textblock*}{3cm}(15cm,23cm) 
\begin{center}
\includegraphics[height=1.2cm]{figs/logouclm.jpg}
\end{center}
\end{textblock*}



\begin{textblock*}{\textwidth}(2.8cm,2.6cm)  % Corrige la posición 
\begin{flushleft}
\includegraphics[height=0.8cm]{figs/logoesii.png} 
\end{flushleft}
\end{textblock*}

\begin{textblock*}{\textwidth}(3cm,2.37cm)  % Corrige la posición
\begin{flushright}
\includegraphics[height=1cm]{figs/logoesi.png} 
\end{flushright}
\end{textblock*}

\begin{textblock*}{\textwidth}(3cm,6.5cm) 
\begin{center}\doublespacing 
{\fontsize{18pt}{4pt}\selectfont \bft{TRABAJO FIN DE MÁSTER}}

\medskip
{\fontsize{18pt}{4pt}\selectfont \espec}


\end{center}
\end{textblock*}


\begin{textblock*}{\textwidth}(3cm,12.5cm) 
\begin{center}\doublespacing 
{\fontsize{22pt}{4pt}\selectfont \bft{\titulo}}
\end{center}
\end{textblock*}


\begin{textblock*}{\textwidth}(3cm, 19.5cm) 
\begin{flushleft}\doublespacing
{\fontsize{14pt}{4pt}\selectfont \bft{Autor:} \autor}

{\fontsize{14pt}{4pt}\selectfont \bft{Director:} \director}
\end{flushleft}
\end{textblock*}


\begin{textblock*}{\textwidth}(3cm,24.75cm) 
\begin{flushright}
{\fontsize{14pt}{4pt}\selectfont \fecha}
\end{flushright}
\end{textblock*}


% -------------------------------
% Dedicatoria
% -------------------------------
\cleardoublepage
\thispagestyle{empty}

\vspace*{9cm}  
\begin{flushright} \em 
Aquí va la dedicatoria que cada cual \\ 
quiera escribir. El ancho se controla \\ 
manualmente
\end{flushright}


% -------------------------------
% Declaración de autoría
% -------------------------------
\cleardoublepage
\thispagestyle{plain}
\chapter*{Declaración de autoría}

%\begin{center}
%\Large{\bft{Declaración de autoría}}
%\end{center}
%\vskip1cm

Yo, \autor, con DNI \quad \dniautor \quad, declaro que soy el único autor del trabajo fin de máster titulado ``\titulo'', que el citado trabajo no infringe las leyes en vigor sobre propiedad intelectual, y que todo el material no original contenido en dicho trabajo está apropiadamente atribuido a sus legítimos autores.

\vspace*{2cm}
\begin{center}
Albacete, a \quad \ldots \quad de \quad \ldots \quad de 20 \ldots

\vskip3cm

Fdo.: \autor
\end{center}


% -------------------------------
% Resumen
% -------------------------------
\cleardoublepage
\thispagestyle{plain}
\chapter*{Resumen}
%\begin{center}
%\Large{\bft{Resumen}}
%\end{center}
%\vskip0.5cm

La información curricular y de trayectoria investigadora se reparte entre
documentos individuales, como los currículos CVN, y registros públicos a
escala masiva, como los perfiles de ORCID. Los sistemas de información de
investigación interoperan datos de este tipo sobre infraestructura
institucional cerrada, y no existe una plataforma abierta, autoalojada y
reproducible que combine ambas naturalezas de dato y las procese de forma
conjunta.

Este Trabajo Fin de Máster diseña, despliega y evalúa una plataforma
\textit{lakehouse} que integra ORCID real con CVN sintético y publica
indicadores de investigación. La plataforma se despliega sobre Kubernetes y
combina almacenamiento de objetos con Iceberg, procesamiento distribuido con
Spark, orquestación con Airflow, publicación en PostgreSQL y visualización
con Superset. Los currículos CVN se generan de forma sintética, válidos
frente al formato Open CVN y sembrados con campos públicos de ORCID, de modo
que no se tratan datos personales de terceros. Esa siembra enlaza cada
currículo con su perfil real y sirve de base a una resolución de identidad
determinista, con evidencia explícita en cada fusión.

La evaluación usa datos reales. Sobre más de 26 millones de registros de
ORCID se filtraron 301\,763 perfiles con afiliación española, y la
resolución de identidad alcanzó una precisión del 95,8\% y una cobertura del
74,2\% en su regla de repliegue por nombre y afiliación. Un banco de pruebas
de rendimiento sobre tres volúmenes de datos demuestra que el número de
ejecutores de Spark es una palanca de fiabilidad, no solo de velocidad: con
la memoria fija por ejecutor, un único ejecutor deja de completar la
transformación a partir del doble del volumen base. Por último, la
plataforma completa se reconstruyó desde cero sobre un clúster aislado
siguiendo únicamente una guía escrita, lo que confirma su reproducibilidad.

% -------------------------------
% Agradecimientos
% -------------------------------
\cleardoublepage
\thispagestyle{plain}
\vspace*{4cm}
\begin{center}
\Large{\bft{Agradecimientos}}
\end{center}
\vskip0.5cm

\lipsum[1]